% 教学页面：五组并列控制变量；连线是作用关系，不表示执行顺序。
\begin{tikzpicture}[
 x=1mm,y=1mm,
 font=\figurefont\small,
 dim/.style={draw=BookBlue,fill=BookBlue!6,line width=0.85pt,
             text width=38mm,align=left,inner sep=3mm},
 card/.style={draw=BookRule,fill=white,line width=0.85pt},
 relation/.style={draw=BookMuted,line width=0.85pt,rounded corners=1.2mm}
]
 \path[use as bounding box] (0,0) rectangle (169,82);
 \draw[draw=BookInk,fill=BookPaper,line width=1.1pt]
      (57,1) rectangle (110,81);
 \node[anchor=west,font=\figurefont\small\bfseries,text=BookInk]
      at (61,76) {通勤与户外};
 \node[anchor=west,font=\figurefont\footnotesize,text=BookMuted]
      at (61,69) {09:40 · 广告已见1次};

 \draw[card] (61,49) rectangle (106,64);
 \node[anchor=west,text=BookInk] at (64,59) {$i_1$　背包测评};
 \node[anchor=west,font=\figurefont\footnotesize,text=BookMuted]
      at (64,53) {重量、收纳与使用条件};

 \draw[draw=BookAmber,fill=BookAmber!7,line width=0.85pt]
      (61,25) rectangle (106,45);
 \node[anchor=west,text=BookInk] at (64,40) {$i_2$　背包商品};
 \node[anchor=east,font=\figurefont\footnotesize,text=BookAmber]
      at (104,40) {广告};
 \node[anchor=west,font=\figurefont\footnotesize,text=BookInk]
      at (64,33) {399元 · 面料防泼水};
 \node[anchor=west,font=\figurefont\footnotesize,text=BookMuted]
      at (64,28) {不适合长时间淋雨};

 \draw[BookMuted,dashed,line width=0.85pt]
      (58,47) -- (109,47);
 \draw[card] (61,5) rectangle (106,21);
 \node[anchor=west,text=BookInk] at (64,16) {$i_3$　保养指南};
 \node[anchor=east,font=\figurefont\footnotesize,text=BookTeal]
      at (103,9) {收藏　不感兴趣};

 \node[dim] (content) at (23,65) {\textbf{内容}\\选入哪些对象、各来源数量};
 \node[dim] (position) at (23,43) {\textbf{位置}\\顺序、模块与首屏范围};
 \node[dim] (expression) at (23,17) {\textbf{表达}\\摘要、条件与商业标识};
 \node[dim] (timing) at (146,68) {\textbf{时间}\\本次时机、间隔与累计频次};
 \node[dim] (interaction) at (146,10) {\textbf{交互}\\可用入口与操作后的响应};

 \draw[relation] (content.east) -- (52,65) -- (52,59) -- (61,59);
 \draw[relation] (position.east) -- (52,43) -- (52,47) -- (58,47);
 \draw[relation] (expression.east) -- (52,17) -- (52,33) -- (61,33);
 \draw[relation] (timing.west) -- (110,68);
 \draw[relation] (interaction.west) -- (110,10);
\end{tikzpicture}
